\chapter[Band 12: Halbordnungen]{Band 12 auf einen Blick}
\noindent\textbf{Halbordnungen}\par\medskip
\noindent\emph{Lesart.} Es gilt \(\PartOrd{A,\leq}\),
\(T\subseteq A\). Ein \emph{minimales} Element wird von keinem
verschiedenen Element unterschritten; ein \emph{Minimum} liegt unter
allen Elementen. Diese Begriffe werden ausdrücklich unterschieden.
\begin{BandUebersicht}
\textbf{Partielle Ordnung und duale Relation}
Für \(x,y,z\in A\):
\UebFormel{\begin{aligned}
&x\leq x,\\
&x\leq y\dsep y\leq z\vdash x\leq z,\\
&x\leq y\dsep y\leq x\vdash x=y,\\
&x\geq y\leftrightarrow y\leq x.
\end{aligned}}
Die ersten drei Zeilen sind Strukturaxiome, die letzte definiert die Dualität.
\UebQuellen{\FormulaRefAuto{Partielle Ordnung}[def];
\FormulaRefAuto{OrderAntisymmetry}[ax];
\FormulaRefAuto{OrderDualRelation}[def]}
&
\textit{Dualität bewahrt die partielle Ordnung}
\UebFormel{\PartOrd{A,\leq}\vdash\PartOrd{A,\geq}.}
\UebQuellen{\FormulaRefAuto{OrderDualPartialOrder}[thm]}\\
\addlinespace[.8ex]
\textbf{Minimales Element und Minimum}
\UebFormel{\begin{aligned}
&\operatorname{MinEl}_{A,\leq}(m,T)\coloneqq\\
&\quad m\in T\land\forall x\in T\,(x\leq m\rightarrow x=m).
\end{aligned}}
Existiert ein kleinstes Element, wird definiert:
\UebFormel{\min_{A,\leq}(T)\coloneqq
\iota m\in T\,\forall x\in T\,(m\leq x).}
\UebQuellen{\FormulaRefAuto{OrderMinimalElement}[def];
\FormulaRefAuto{Minimum}[def]}
&
\textit{Eindeutigkeit und Minimalität des Minimums}
\UebFormel{\begin{aligned}
&\exists m\in T\,\forall x\in T\,(m\leq x)\\
&\quad\vdash\exists!m\in T\,\forall x\in T\,(m\leq x).
\end{aligned}}
Unter dieser Existenzvoraussetzung gilt außerdem
\UebFormel{\operatorname{MinEl}_{A,\leq}(\min_{A,\leq}(T),T).}
\UebQuellen{\FormulaRefAuto{MinimumUnique}[thm];
\FormulaRefAuto{MinimumIsMinimalElement}[thm]}\\
\addlinespace[.8ex]
\textbf{Maximales Element und Maximum}
\UebFormel{\begin{aligned}
&\operatorname{MaxEl}_{A,\leq}(m,T)\coloneqq\\
&\quad m\in T\land\forall x\in T\,(m\leq x\rightarrow x=m).
\end{aligned}}
Existiert ein größtes Element, wird definiert:
\UebFormel{\max_{A,\leq}(T)\coloneqq
\iota m\in T\,\forall x\in T\,(x\leq m).}
\UebQuellen{\FormulaRefAuto{OrderMaximalElement}[def];
\FormulaRefAuto{Maximum}[def]}
&
\textit{Eindeutigkeit und Maximalität des Maximums}
\UebFormel{\begin{aligned}
&\exists m\in T\,\forall x\in T\,(x\leq m)\\
&\quad\vdash\exists!m\in T\,\forall x\in T\,(x\leq m).
\end{aligned}}
Unter dieser Existenzvoraussetzung gilt
\UebFormel{\operatorname{MaxEl}_{A,\leq}(\max_{A,\leq}(T),T).}
\UebQuellen{\FormulaRefAuto{MaximumUnique}[thm];
\FormulaRefAuto{MaximumIsMaximalElement}[thm]}\\
\addlinespace[.8ex]
\textbf{Restriktion einer Ordnung}
Für \(x,y\in T\):
\UebFormel{x\leq_T y\leftrightarrow x\leq y.}
\UebQuellen{\FormulaRefAuto{OrderRestrictionOperator}[def]}
&
\textit{Die Teilmenge erbt eine partielle Ordnung}
\UebFormel{\begin{aligned}
&T\subseteq A\dsep\PartOrd{A,\leq}\\
&\hspace{12mm}\vdash\PartOrd{T,\leq_T}.
\end{aligned}}
\UebQuellen{\FormulaRefAuto{OrderRestrictionPartialOrder}[thm]}\\
\addlinespace[.8ex]
\textbf{Hauptfilter}
\UebFormel{\PrincipalFilter{A,\leq}{j}\coloneqq
\{x\in A\mid j\leq x\}.}
\UebQuellen{\FormulaRefAuto{PrincipalFilter}[def]}
&
\textit{Typisierung}
\UebFormel{\PrincipalFilter{A,\leq}{j}\subseteq A.}
Der Hauptfilter sammelt genau die Elemente oberhalb von \(j\).
\UebQuellen{\FormulaRefAuto{PrincipalFilterSubsetCarrier}[thm]}\\
\end{BandUebersicht}
